\section{Introduction}

Water-quality monitoring in remote Amazonian areas is constrained by sparse infrastructure, high field logistics cost, and intermittent network connectivity. Edge AI deployment is the central motivation of this work: local, fully offline inference is required to keep monitoring functional under unstable communications and limited energy budgets. In this context, continuous dependence on cloud inference is operationally fragile and energetically expensive, especially for battery-powered deployments and long-duration campaigns. At the same time, water-quality variables such as pH, dissolved oxygen (OD), and turbidity are highly relevant for environmental screening and public-health risk mitigation.

Conventional machine-learning pipelines can provide strong predictive performance, but many are designed for server-centric execution and do not directly address energy constraints at the edge. Neuromorphic models based on spiking neural networks (SNNs) offer an event-driven alternative compatible with low-power embedded systems, yet practical evidence in real environmental sensing workflows remains limited, particularly when combining multimodal prediction and platform-level energy benchmarking.

This paper proposes and evaluates an offline multimodal SNN pipeline for water-quality prediction in an Amazonian field scenario, including end-to-end deployment and cross-platform energy analysis. The study integrates data processing, model training, geospatial interpretation, and embedded inference on heterogeneous hardware (ESP32, Raspberry Pi, and Android devices).

The main contributions of this work are \rev{fourfold}: (i) the development of a multimodal SNN-based water-quality prediction pipeline designed for offline edge deployment; (ii) the creation of an augmented georeferenced dataset derived from field measurements in Lago do Serpa; (iii) a comparative latency--energy evaluation across heterogeneous platforms, including ESP32, Raspberry Pi, and Android devices\rev{; and (iv) a statistically grounded evaluation protocol for small-data environmental regression, combining repeated-seed benchmarking, confidence intervals, and nonparametric significance testing. Beyond energy benchmarking, the scientific contribution lies in demonstrating, end to end, that multimodal neuromorphic regression of physicochemical water-quality targets is feasible from field acquisition to fully offline inference on three distinct classes of edge hardware, a combination not yet reported for freshwater monitoring in the Amazon region.}
